% Compile twice: lualatex 30005052.tex
\documentclass[11pt]{article}
\usepackage[margin=25mm]{geometry}
\usepackage{fontspec,unicode-math}
\setmainfont{Latin Modern Roman}
\setmathfont{Latin Modern Math}
\usepackage{amsmath,amssymb,mathtools}
\usepackage{xurl,hyperref}
\hypersetup{colorlinks=true,urlcolor=blue,linkcolor=blue}
\setcounter{secnumdepth}{0}
\setlength{\parindent}{0pt}
\setlength{\parskip}{5pt}
\setlength{\emergencystretch}{3em}
\newcommand{\sref}[2]{\hyperlink{source:#1:#2}{[S#2]}}
\newcommand{\cataloguescope}[1]{\paragraph{Source question.}#1}
\newcommand{\floor}[1]{\left\lfloor#1\right\rfloor}
\begin{document}
% UnsolvedMath ID 30005052; source catalogue code OWR-9790365-001
\setcounter{section}{30005051}
\section{Local-to-Global Tensorization of Information Functionals}\label{30005052}
{\small\sffamily UnsolvedMath ID 30005052 · Computer Science}
\subsection{Definitions and mathematical statement}

\paragraph{Shared notation.}Standard number systems are denoted by $\mathbb N,\mathbb Z,\mathbb Q,\mathbb R,\mathbb C$; index ranges, computational models and specialized notation follow the cited problem source and the definitions below.


For finite-alphabet random variables, $H(Z)=-\sum_z p(z)\log p(z)$ is entropy, with $0\log0=0$. In the source family a channel $W$ sends $X$ to $(Y_1,\ldots,Y_k)$, and real weights $\alpha_S$ and an input penalty $\gamma$ specify
\[F_{\gamma,W}(p_X)=\sum_{S\subseteq[k]}\alpha_S H(Y_S)-\mathbb E[\gamma(X)].\]
For two independent channels, the joint functional replaces $H(Y_S)$ by $H(Y_{1,S},Y_{2,S})$ and subtracts $\mathbb E[\gamma_1(X_1)+\gamma_2(X_2)]$, allowing correlated inputs. Local tensorization means products of local maximizers of the separate functionals are local maximizers of the joint functional, for every pair of penalties. Global tensorization means a product input attains the joint maximum for every pair of penalties.
\paragraph{Statement recorded in the supplied source.}
For which families of information functionals does local tensorization imply global tensorization? \sref{30005052}{2}
\paragraph{Scope qualification.}
The functional and the local/global distinction are recovered from the report discussion on pp. 689--690. Local maximality is relative to the finite probability simplex; the question asks for families, rather than asserting an implication for all information quantities.
\paragraph{Source status.}
The supplied catalogue labels this record open. Its source-status metadata was checked against the supplied snapshot on 2026-10-11. This import does not independently establish current open status or mathematical correctness. The catalogue literature assessment is dated 2026-08-22.
\subsection{Short English statement}
For which channel and entropy-weight families does the local product-maximizer property force a global product-maximizer property?
\subsection{Sources}
\begin{itemize}
\item[\hypertarget{source:30005052:1}{S1}] Ulam AI and the UnsolvedMath Contributors, UnsolvedMath: A Curated Collection of Open Mathematics Problems, record 30005052 (OWR-9790365-001). Catalogue curation and metadata: CC BY 4.0. The supplied source question is retained; definitions and the English summary are editorial additions.
\url{https://huggingface.co/datasets/ulamai/UnsolvedMath}.
\item[\hypertarget{source:30005052:2}{S2}] New Mathematical Techniques in Information Theory (2023), Oberwolfach Reports, DOI 10.4171/owr/2022/14, pp. 683--707
\url{https://doi.org/10.4171/owr/2022/14}.
\item[\hypertarget{source:30005052:3}{S3}] Nair, local and global tensorization
\url{https://ems.press/content/serial-article-files/46950?nt=1}.
\item[\hypertarget{source:30005052:4}{S4}] New Mathematical Techniques in Information Theory, Oberwolfach Report 14/2022, pp. 689--690.
\url{https://ems.press/content/serial-article-files/46950}.
\textit{Underlying Oberwolfach Reports material retains its source-specific licence, including CC BY-SA 4.0 where applicable; see THIRD\_PARTY\_NOTICES.md.}
\end{itemize}
\paragraph{Preparation and review.}
Original question attributed to the cited source and UnsolvedMath contributors. Codex (GPT-6; exact serving revision not exposed) assisted with definitions, English summaries, source inspection and formatting. No new human mathematical authorship or independent proof verification is asserted. Canonical ID 30005052 and source code OWR-9790365-001 checked against the complete supplied registry and existing repository statements; known flagged repetitions and obvious identical conjectures excluded. Full mathematical-equivalence/literature review has not been performed. Statement preparation and final review followed contribution prompts C and D; automated checks verify structure and identity, not open status.
\label{end:30005052}
\end{document}
